% ============================================================
%  figures_and_tables_guide.tex
%  Reference sheet for inserting figures and tables in LaTeX.
%  Keep this in thesis/ — it is NOT \input'd into main.tex,
%  just a reference document for you to consult while writing.
% ============================================================

% ==============================================================
%  PART 1: FIGURES
% ==============================================================

% ---- Basic figure -------------------------------------------
% Save all figures to thesis/figures/ as .pdf or .png.
% .pdf is preferred for vector graphics (plots from matplotlib).
% .png is fine for raster images.

\begin{figure}[H]                          % [H] = "place HERE exactly"
    \centering
    \includegraphics[width=0.85\textwidth]{figures/my_figure.pdf}
    \caption{A descriptive caption explaining what the figure shows.}
    \label{fig:my_figure}                  % Used for \ref{fig:my_figure}
\end{figure}

% Width options:
%   width=0.85\textwidth    — 85% of text width (most common)
%   width=0.5\textwidth     — half-width
%   width=\textwidth        — full width
%   height=3in              — fixed height, aspect ratio preserved

% ---- Two figures side by side --------------------------------
\begin{figure}[H]
    \centering
    \begin{subfigure}[b]{0.48\textwidth}
        \centering
        \includegraphics[width=\textwidth]{figures/fig_left.pdf}
        \caption{Left panel caption.}
        \label{fig:left}
    \end{subfigure}
    \hfill
    \begin{subfigure}[b]{0.48\textwidth}
        \centering
        \includegraphics[width=\textwidth]{figures/fig_right.pdf}
        \caption{Right panel caption.}
        \label{fig:right}
    \end{subfigure}
    \caption{Overall figure caption (shown below both panels).}
    \label{fig:both}
\end{figure}

% ---- Saving figures from Python for LaTeX -------------------
% In your Python script, save as PDF for best quality:
%
%   import matplotlib.pyplot as plt
%   fig, ax = plt.subplots(figsize=(7, 4))
%   ax.plot(x, y)
%   ax.set_xlabel("X axis")
%   ax.set_ylabel("Y axis")
%   plt.tight_layout()
%   plt.savefig("thesis/figures/my_figure.pdf", dpi=300, bbox_inches="tight")
%
% OR save as high-res PNG:
%   plt.savefig("thesis/figures/my_figure.png", dpi=300, bbox_inches="tight")


% ==============================================================
%  PART 2: TABLES
% ==============================================================

% ---- Basic table with booktabs (econometrics style) ----------
% booktabs gives you \toprule, \midrule, \bottomrule
% NEVER use \hline in economics tables — use booktabs rules

\begin{table}[H]
    \centering
    \caption{Table title goes here (above the table in economics style).}
    \label{tab:basic}
    \begin{tabular}{lcc}        % l = left, c = center, r = right
        \toprule
        Variable & Column 1 & Column 2 \\
        \midrule
        Row 1    & 0.123    & 0.456    \\
        Row 2    & 0.789$^{**}$ & 0.012$^{*}$ \\
                 & (0.034)  & (0.007)  \\       % std errors in parens below
        \midrule
        N        & 1,000    & 1,000    \\
        R$^2$    & 0.45     & 0.52     \\
        \bottomrule
    \end{tabular}
\end{table}

% ---- Significance stars convention --------------------------
% Standard in economics:
%   $^{*}$    p < 0.10
%   $^{**}$   p < 0.05
%   $^{***}$  p < 0.01
% Report in table notes, not in column headers.

% ---- Table with notes (threeparttable) ----------------------
% Use threeparttable when you need footnotes below the table.
% This is the standard for regression tables.

\begin{table}[H]
    \centering
    \caption{Regression Results with Notes}
    \label{tab:with_notes}
    \begin{threeparttable}
        \begin{tabular}{lcc}
            \toprule
            & (1) OLS & (2) OLS + FE \\
            \midrule
            Violence (early childhood) & $-$0.022$^{***}$ & 0.010 \\
                                       & (0.007)          & (0.010) \\[0.3em]
            Controls & Yes & Yes \\
            LGA FE   & No  & Yes \\
            \midrule
            N        & 56,878 & 56,878 \\
            R$^2$    & 0.43   & 0.52   \\
            \bottomrule
        \end{tabular}
        \begin{tablenotes}
            \small
            \item \textit{Notes:} Robust standard errors in parentheses.
            $^{*}p < 0.10$, $^{**}p < 0.05$, $^{***}p < 0.01$.
        \end{tablenotes}
    \end{threeparttable}
\end{table}

% ---- Aligning decimal points (siunitx) ----------------------
% Use S column type to align numbers at the decimal point.
% Requires \usepackage{siunitx} in preamble.

\begin{table}[H]
    \centering
    \caption{Decimal-Aligned Table}
    \label{tab:siunitx}
    \begin{tabular}{l S[table-format=1.3] S[table-format=1.3]}
        \toprule
        Variable & {Column 1} & {Column 2} \\   % braces prevent formatting
        \midrule
        Coef 1   & 0.123 & -0.456 \\
        Coef 2   & 1.789 & 0.012  \\
        \bottomrule
    \end{tabular}
\end{table}

% ---- Wide table in landscape (pdflscape) --------------------
% For tables too wide to fit in portrait. Requires \usepackage{pdflscape}.

\begin{landscape}
\begin{table}[H]
    \centering
    \caption{Wide Table in Landscape}
    \label{tab:landscape}
    \begin{tabular}{lcccccc}
        \toprule
        Variable & Col1 & Col2 & Col3 & Col4 & Col5 & Col6 \\
        \midrule
        Row 1    & & & & & & \\
        \bottomrule
    \end{tabular}
\end{table}
\end{landscape}

% ---- Multi-page table (longtable) ---------------------------
% For variable definition tables or long appendix tables.

\begin{longtable}{p{0.25\textwidth}p{0.65\textwidth}}
    \caption{Long Table with Page Breaks} \label{tab:long} \\
    \toprule
    Variable & Definition \\
    \midrule
    \endfirsthead                           % Everything above repeats on next pages:
    \multicolumn{2}{c}{\tablename~\thetable{} (continued)} \\
    \toprule
    Variable & Definition \\
    \midrule
    \endhead
    \midrule
    \multicolumn{2}{r}{\textit{Continued on next page}} \\
    \endfoot
    \bottomrule
    \endlastfoot
    % Table content here:
    var1 & Definition of variable 1. \\
    var2 & Definition of variable 2. \\
\end{longtable}

% ---- Generating LaTeX tables from Python --------------------
% Option A: use the `stargazer` equivalent in Python —
%   pip install stargazer
%   from stargazer.stargazer import Stargazer
%   import statsmodels.formula.api as smf
%   model = smf.ols("attend ~ viol_ech + age + male", data=df).fit()
%   s = Stargazer([model])
%   print(s.render_latex())   # paste output into your .tex file
%
% Option B: use pandas DataFrame.to_latex()
%   df.to_latex("thesis/tables/summary_stats.tex", index=False)
%   Then \input{tables/summary_stats.tex} in your chapter file.
%
% Option C: modelsummary (if using R)
%   library(modelsummary)
%   modelsummary(model, output = "thesis/tables/results.tex")


% ==============================================================
%  PART 3: CROSS-REFERENCING
% ==============================================================

% Reference a figure:   See Figure~\ref{fig:my_figure}.
% Reference a table:    See Table~\ref{tab:basic}.
% Reference a section:  See Section~\ref{sec:methodology}.
% Reference an equation: See Equation~\eqref{eq:attend}.

% The tilde (~) prevents a line break between "Figure" and the number.
% Always use \eqref for equations (adds parentheses automatically).

% ==============================================================
%  PART 4: WORKFLOW SUMMARY
% ==============================================================

% 1. Write Python script → save figure to thesis/figures/my_fig.pdf
% 2. In your .tex chapter file:
%       \begin{figure}[H]
%           \centering
%           \includegraphics[width=0.85\textwidth]{figures/my_fig.pdf}
%           \caption{Caption.}
%           \label{fig:my_fig}
%       \end{figure}
% 3. Reference in text: "...as shown in Figure~\ref{fig:my_fig}."
% 4. Compile with pdflatex (Overleaf does this automatically).

% For tables from Python:
% 1. Export .tex snippet using stargazer or to_latex()
% 2. Save to thesis/tables/table_name.tex
% 3. In your chapter file: \input{tables/table_name.tex}
%    OR paste the tabular environment directly into the chapter.
